\section{Discussion and Limitations}
Shared anchors solve a coordinate-consistency problem: pose corrections carry
local maps and their observations together. The experiments show that this
design supports appearance refinement with the inherited tracker. They also
show where dense prediction needs further work.

First, overlapping predictions create expensive maps. Removing primitives
after optimisation loses useful coverage. A more compact system should
allocate and refine primitives under a fixed budget from the start.
A second GPU reduces tracker contention, but the measured throughput is
still below 30\,Hz and the best reconstruction results include 300\,s polish.

Second, shared anchors preserve local relationships without merging
disagreeing surfaces. Seams between anchors and colour differences between
source images can remain. Joint treatment of pose, geometry, and appearance
may help reduce these errors.

Finally, the external comparison covers eight Replica scenes and one TUM
scene, with different optimisation budgets and map sizes across systems.
A broader study should use repeated runs, equal budgets, and scoring views
excluded from every system's optimisation. The near-zero adaptation and
confidence-weighting effects especially need replication.

\section{Conclusion}
Splatt3R-SLAM turns image-pair predictions into a persistent map by anchoring
Gaussians and supervision cameras together. This lets geometric tracking
and appearance refinement proceed separately while sharing the same pose
graph. Common-camera evaluation shows a 3.54\,dB Replica PSNR gain over
Photo-SLAM and close PSNR to MonoGS on TUM desk. The next step is to retain
this quality with fewer primitives and a bounded online refinement budget.
